% =============================================================
%  Section: Shipai Full-Scenario Benchmark Results
%  Auto-generated from shipai_strategy_analysis.py
%  Requires packages: graphicx, booktabs, amsmath, siunitx
% =============================================================

\subsection{Shipai Full-Scenario Benchmark}
\label{sec:shipai}

We evaluate the three navigation strategies---Static, Density-aware
($\Delta t_u = 60\,\mathrm{s}$), and Density-Hazard
($\beta = 1,\,\Delta t_u = 60\,\mathrm{s}$)---on the Shipai
district full-scenario benchmark.
Each configuration is evaluated over $N = 50$ independent runs
with randomised initial conditions; all reported values are means with
95\,\% bootstrap confidence intervals (Table~\ref{tab:shipai_comparison}).

% -------------------------------------------------------------------
\subsubsection{Evacuation Efficiency}
\label{sec:shipai_efficiency}
% -------------------------------------------------------------------

The static baseline completes 95\,\% evacuation at
$T_{0.95} = 1114\,\mathrm{s}$
($T_{0.99} = 1431\,\mathrm{s}$).
Both dynamic strategies substantially accelerate crowd clearance
(Figure~\ref{fig:shipai_bar4}, panels a--b).
The Density-aware strategy reduces $T_{0.95}$ by
14.1\,\% to $956\,\mathrm{s}$
and $T_{0.99}$ by 14.2\,\% to
$1227\,\mathrm{s}$.
The Density-Hazard strategy achieves a comparable reduction in
$T_{0.95}$ (13.7\,\%, to $961\,\mathrm{s}$)
and a marginally larger improvement in $T_{0.99}$
(16.8\,\%, to $1191\,\mathrm{s}$),
reflecting a tighter tail of the evacuation-time distribution.
The individual evacuation-time distributions (Figure~\ref{fig:shipai_violin})
confirm that both dynamic strategies shift the entire distribution
leftward, with the Density-Hazard variant producing the narrowest
99th-percentile.

% -------------------------------------------------------------------
\subsubsection{Risk Reduction}
\label{sec:shipai_risk}
% -------------------------------------------------------------------

The most pronounced differentiation between the two dynamic strategies
emerges in the risk metrics (Figure~\ref{fig:shipai_bar4}, panels c--d;
Figure~\ref{fig:shipai_scatter}).
Starting from a static-baseline cumulative soft risk of
$\Sigma_s = 4.31e+07\,\mathrm{person\cdot s}$,
the Density-aware strategy achieves a moderate reduction of
38.3\,\% (to $2.66e+07$).
The Density-Hazard strategy delivers a far larger reduction of
78.2\,\%\,, collapsing $\Sigma_s$ to
$9.39e+06\,\mathrm{person\cdot s}$---approximately
$4.6\times$ lower than
the static baseline and
$2.8\times$ lower than
the Density-aware variant.
This gap highlights that routing agents away from high-density cells
(Density-aware) and routing them away from high-hazard cells
(Density-Hazard) are complementary but distinct objectives; the latter
is substantially more effective at limiting cumulative exposure to
dangerous zones.

The hard-harm metric $H$ (the number of agents simultaneously exposed
to hazard above the critical threshold) mirrors this pattern.
The static baseline incurs
$H = 2.65e+06$
(95\,\% CI
[1.61e+06, 3.74e+06]).
The Density-aware strategy reduces $H$ by 83.0\,\% to
$4.51e+05$,
whereas the Density-Hazard strategy nearly eliminates hard harm,
achieving $H = 2863.9$
(99.9\,\% reduction).

% -------------------------------------------------------------------
\subsubsection{Spatial Congestion Patterns}
\label{sec:shipai_spatial}
% -------------------------------------------------------------------

The spatial-efficiency analysis (Figure~\ref{fig:shipai_spatial})
quantifies how crowd energy is distributed across the simulation domain.
The mean density-dwell time (layer-7 mean, the average duration a cell
is occupied above the minimum threshold) decreases from
$18.88\,\mathrm{s/cell}$ (Static) to
$17.26\,\mathrm{s/cell}$ (Density-aware) and
$17.24\,\mathrm{s/cell}$ (Density-Hazard),
a reduction of
8.6\,\% and
8.7\,\%,
respectively.
The risk-weighted dwell time (layer-8 mean) follows the same ordering
but with a larger absolute reduction for the Density-Hazard strategy
(from $1.800$ to $1.465\,\mathrm{s/cell}$,
a 18.6\,\%
improvement), consistent with the risk-weighted routing objective.
The high-congestion area (cells with dwell time $\geq$ Q95) contracts
modestly under both dynamic strategies, indicating that density-aware
routing disperses persistent hot-spots without significantly altering
the global congestion footprint.

% -------------------------------------------------------------------
\subsubsection{Agent Behaviour}
\label{sec:shipai_behavior}
% -------------------------------------------------------------------

Figure~\ref{fig:shipai_behavior} examines the behavioural
consequences of dynamic routing at the individual level.
Mean travel distance increases slightly from
$137.6\,\mathrm{m}$ (Static) to
$145.6\,\mathrm{m}$ (Density-aware) and
$144.9\,\mathrm{m}$ (Density-Hazard),
indicating that dynamic guidance induces modest path lengthening
in exchange for improved global efficiency.
The mean route-switching count rises from
$0.069$ (Static) to
$0.143$ (Density-aware) and
$0.122$ (Density-Hazard), reflecting the
additional guidance-field updates but remaining below one switch
per agent on average.

% -------------------------------------------------------------------
\subsubsection{Computational Overhead}
\label{sec:shipai_compute}
% -------------------------------------------------------------------

Both dynamic strategies incur substantially higher computation costs
than the static baseline, owing to repeated dynamic-field recomputation
(Figure~\ref{fig:shipai_compute}).
The Density-aware strategy requires
$3.2\times$ the static
total computation time
($337\,\mathrm{s}$ vs.\
$106\,\mathrm{s}$),
of which $71.9\,\%$
is attributable to field updates
($242\,\mathrm{s}$).
The Density-Hazard strategy shows a slightly lower total cost
($337\,\mathrm{s}$,
$3.2\times$ baseline),
with $72.4\,\%$
spent on field updates.
This overhead is commensurate with the $13.7\,\% reduction
in $T_{0.95}$ and the dramatic risk reduction achieved.

% -------------------------------------------------------------------
\subsubsection{Summary}
\label{sec:shipai_summary}
% -------------------------------------------------------------------

Taken together, the Shipai benchmark results demonstrate that the
Density-Hazard strategy ($\beta = 1,\,\Delta t_u = 60\,\mathrm{s}$)
Pareto-dominates the Density-aware baseline across all risk metrics
while achieving comparable or superior evacuation efficiency.
The $78.2\,\% reduction in cumulative soft risk
and the near-elimination of hard harm
confirm that incorporating hazard information into the routing
potential is the key differentiator, rather than the density-awareness
alone.
The moderate increase in computation cost
($3.2\times$ baseline)
is well within the budget established by the update-period sensitivity
analysis (Section~\ref{sec:update_period}), validating
$\Delta t_u^{*} = 60\,\mathrm{s}$ as the operating point.


\begin{table}[htbp]
\centering
\caption{%
  Performance comparison of three navigation strategies on the Shipai
  full-scenario benchmark ($N = 50$ independent runs each).
  Bracketed ranges are 95\,\% bootstrap confidence intervals.
  $T_{0.95}$/$T_{0.99}$: time for 95/99\,\% of agents to evacuate (s).
  $\Sigma_s$: cumulative soft risk (person$\cdot$s).
  $H$: hard-harm count.
  AllTime: mean total computation time (ms).%
}
\label{tab:shipai_comparison}
\begin{tabular}{@{}lccccc@{}}
\toprule
\textbf{Strategy} & $T_{0.95}$ (s) & $T_{0.99}$ (s) & $\Sigma_s$ & $H$ & AllTime (ms) \\
\midrule
Static & 1.1k [1.1k, 1.2k] & 1.4k [1.4k, 1.5k] & 43.11M [32.84M, 54.24M] & 2.65M [1.61M, 3.74M] & 106.1k \\
Density-60s & 956.2 [930.3, 982.4] & 1.2k [1.2k, 1.3k] & 26.59M [22.19M, 31.24M] & 451.5k [267.0k, 654.4k] & 337.0k \\
D-H $\beta=1$, 60s & 960.9 [933.4, 987.7] & 1.2k [1.1k, 1.2k] & 9.39M [8.64M, 10.25M] & 2.9k [670, 5.9k] & 336.8k \\
\bottomrule
\end{tabular}
\end{table}

